\section{The greatest ambit}\label{sec:ambit}

Every topological group has a largest flow with a distinguished point of dense
orbit. It is the Samuel compactification of the group, and it is where the
universal minimal flow lives.

\begin{definition}[Ambit]\label{def:ambit}
\uses{def:flow}
A \emph{$G$-ambit} is a pair $(X,x_0)$ consisting of a $G$-flow $X$ and a point
$x_0\in X$ whose orbit $Gx_0$ is dense in $X$. A \emph{morphism of ambits}
$(X,x_0)\to(Y,y_0)$ is a $G$-map sending $x_0$ to $y_0$. There is at most one
such morphism, since $Gx_0$ is dense, and it is surjective.
\end{definition}

\begin{definition}[Right uniformly continuous functions]\label{def:ruc}
A bounded function $f\colon G\to\mathbb{C}$ is \emph{right uniformly
continuous} if for every $\varepsilon>0$ there is a neighbourhood $V$ of the
identity such that $|f(vg)-f(g)|<\varepsilon$ for all $g\in G$ and $v\in V$.
These functions form a unital C*-subalgebra $\mathrm{RUC}_b(G)$ of
$\ell^\infty(G)$.
\end{definition}

\begin{remark}[Conventions]\label{rem:ruc-conventions}
\uses{def:ruc}
These are the functions that are uniformly continuous for the uniform
structure with basic entourages $\{(x,y) : xy^{-1}\in V\}$. This structure is
called the right uniformity by Bourbaki and Pestov and the left uniformity by
Ellis and several other authors, so the same algebra appears in the literature
under both names. Check each source's convention before comparing statements.
\end{remark}

\begin{definition}[Samuel compactification]\label{def:samuel}
\uses{def:ruc, def:ambit}
The \emph{Samuel compactification} $S(G)$ of $G$ is the Gelfand spectrum of
$\mathrm{RUC}_b(G)$: the space of characters of $\mathrm{RUC}_b(G)$ with the
weak* topology. For $f\in\mathrm{RUC}_b(G)$ and $h\in G$ the function
$f^h(g)=f(hg)$ is again right uniformly continuous, and
$(h\varphi)(f)=\varphi(f^h)$ defines a continuous action of $G$ on $S(G)$.
Evaluation at points embeds $G$ equivariantly as a dense subset of $S(G)$, so
$(S(G),1)$ is a $G$-ambit.
\end{definition}

\begin{proposition}[Universal property]\label{prop:greatest-ambit}
\uses{def:samuel}
For every $G$-ambit $(X,x_0)$ there is a unique morphism of ambits
$(S(G),1)\to(X,x_0)$. For this reason $S(G)$ is called the \emph{greatest
ambit} of $G$.
\end{proposition}

\begin{proof}[Sketch]
Continuity of the action and compactness of $X$ give, for every
$f\in C(X)$ and $\varepsilon>0$, a neighbourhood $V$ of the identity with
$|f(vx)-f(x)|<\varepsilon$ for all $x\in X$ and $v\in V$. Hence
$g\mapsto f(gx_0)$ lies in $\mathrm{RUC}_b(G)$. This defines a unital
*-homomorphism $C(X)\to\mathrm{RUC}_b(G)=C(S(G))$, whose dual map
$S(G)\to X$ is the required morphism. Uniqueness holds because $G$ is dense in
$S(G)$. See \cite{Auslander1988,Pestov2006}.
\end{proof}
